\section{A sharp exchange budget and constructive repair of repeated seeds}
\label{sec:completion-budget}

The classical unextendible $4\times12$ seed does admit a constructive
repair: deleting one of its four rows permits a full completion.
Tensor copies of this repair give a useful distinction between two
costs. A single augmenting exchange can remain cheap even when the
minimum number of original rows that must be discarded for a full
completion grows without bound. The construction underlying the repair
is classical Scarpis machinery, not a new Hadamard existence theorem;
see Djokovi\'c, \emph{Generalization of Scarpis's theorem on Hadamard
matrices}, \url{https://arxiv.org/abs/1601.00635}, especially its
construction retaining all but one repeated input row. We give a
self-contained exchange-budget argument and the order-twelve repair.
No bibliographic priority is asserted for the budget interpretation.

\begin{theorem}[Odd-block exchange budget]
Let $K$ be a Hadamard matrix of order $m$, let $s>1$ be odd, and set
$P=K\otimes(1,\ldots,1)$, with $s$ repeated entries in each block.
Suppose $m-k$ original rows are retained, and $L$ additional sign rows
are inserted so that all retained and inserted rows are mutually
orthogonal. Then
\[
 L\le k s^2.
\]
In particular, completing a retained subset to an order-$ms$ Hadamard
matrix requires deleting at least
\[
 k\ge\left\lceil\frac{m}{s+1}\right\rceil
\]
original rows. The conclusion allows arbitrary replacement rows and
does not assume product structure in the inserted rows.
\end{theorem}
\begin{proof}
For a sign row $v$ let $b_j(v)$ be its sum on block $j$ and put
$E(v)=\sum_{j=1}^m b_j(v)^2$. Every block has odd size, so $E(v)\ge m$.
An original repeated row has $E(v)=ms^2$.
Let $e_j$ be the indicator of block $j$. Orthogonal sign rows have
squared norm $ms$, so Bessel's inequality gives
\[
 \sum_v\langle v,e_j\rangle^2\le ms\|e_j\|^2=ms^2.
\]
Summing over the $m$ blocks and separating original and inserted rows,
\[
 (m-k)ms^2+Lm\le\sum_v E(v)\le m^2s^2.
\]
Thus $L\le ks^2$. A full completion has
$L=ms-(m-k)=m(s-1)+k$, so
$m(s-1)\le k(s^2-1)$, proving the deletion bound.
\end{proof}

The case $k=0$ recovers unextendibility. Equality $L=ks^2$ also forces
every inserted row to have block sums in $\{+1,-1\}$ and equality in
each blockwise Bessel inequality. Thus the bound identifies a concrete
structure that a maximal-efficiency repair must exploit.

\subsection{The order-twelve repair and its exact candidate code}
Take $K=H_4$, with first row $h=(1,1,1,1)$ and three remaining
orthogonal sign rows. Delete the repeated first row from
$P=H_4\otimes(1,1,1)$. Orthogonality to the other three rows means
the four block sums of a candidate row form $\lambda h$.
Each block sum belongs to $\{-3,-1,1,3\}$, so up to negating the
candidate, $\lambda$ is $1$ or $3$.

The case $\lambda=3$ is exactly the deleted row. For $\lambda=1$,
each block contains one negative entry and two positive entries.
Encode the negative positions by a word in $\mathbb F_3^4$.
This yields exactly $3^4=81$ candidate classes, besides the deleted
row, hence $82$ antipodal candidate classes in total.
Two encoded rows have inner product $4a-4$, where $a$ is the number
of coordinates in which their words agree. Therefore they are
orthogonal exactly when their words agree once. The deleted row
has inner product $4$ with every encoded row, so cannot be retained
alongside any of them.

The nine words
\[
 c(x,y)=(x,y,x+y,x-y),\qquad x,y\in\mathbb F_3,
\]
agree pairwise in exactly one coordinate. To see this, subtract two
words. If one of the two differences is zero, exactly its coordinate
vanishes; otherwise their quotient is $1$ or $-1$, so exactly one
of their sum and difference vanishes. The corresponding nine sign
rows, together with the three retained seed rows, form a Hadamard
matrix $B$ of order $12$.

Conversely, every full completion of the three retained rows must
use nine encoded rows. Indeed the exchange budget is saturated with
$k=1,L=9$, forcing absolute block sums $1$. Its nine words agree
pairwise once. Projection of these words onto any two coordinates is
injective and hence bijective onto $\mathbb F_3^2$; they form an
orthogonal array $\operatorname{OA}(9,4,3,2)$.
Thus the finite affine-plane code is both a sufficient repair and a
necessary combinatorial description of every completion in this
normalization. Deleting any of the four original rows gives the same
description after block sign changes.

\begin{theorem}[Cheap augmentation versus exact linear completion cost]
Let $T$ be a Hadamard matrix of order $t$ and set
\[
 P_t=(H_4\otimes(1,1,1))\otimes T.
\]
This is a $4t\times12t$ partial Hadamard matrix. It has the following
properties, allowing a common column permutation throughout.
\begin{enumerate}
\item It has no additional orthogonal sign row.
\item Its minimum deletion count for an augmenting exchange is exactly
one: one old row can be replaced by two orthogonal new rows.
\item Its minimum deletion count permitting full completion is exactly
$t$.
\item For every $0\le k\le t$, the exchange budget $L\le9k$ is
attained: one can delete $k$ old rows and insert $9k$ new rows.
Every such intermediate matrix is unextendible by an additional row.
\end{enumerate}
Consequently these two notions of repair cost separate arbitrarily
far along, for example, the Sylvester orders $t=2^a$.
\end{theorem}
\begin{proof}
Permuting tensor factors identifies the seed with
$(H_4\otimes T)\otimes(1,1,1)$, so the exchange-budget theorem gives
unextendibility and a full-completion lower bound
$\lceil4t/4\rceil=t$.
For each row $\tau$ of $T$, the four seed rows with second factor
$\tau$ form a fiber. In that fiber replace $h\otimes(1,1,1)\otimes\tau$
by the nine encoded rows of the order-twelve repair, tensored with
$\tau$. Different fibers remain orthogonal because their $T$ rows
are orthogonal. Repairing any $k$ fibers deletes $k$ old rows and
inserts $9k$ rows, proving the fourth assertion. Repairing all fibers
produces $B\otimes T$, a full completion after exactly $t$ deletions.
An additional row at any intermediate stage would exceed the bound
$L\le9k$, so every intermediate matrix is unextendible.
For a single augmenting exchange use just two of the nine available
rows in one fiber; zero deletions are impossible by unextendibility.
\end{proof}

In particular, taking $k=t-1$ gives an explicit unextendible
$(12t-8)\times12t$ partial Hadamard matrix for every Hadamard order
$t$. Thus the obstruction at exactly eight missing rows persists at
unbounded orders, while one further fiber repair still completes it.

This family does not obstruct an iterative constructive strategy:
the displayed sequence of fiber repairs succeeds explicitly. It does
show that a uniform bound on the cost of one augmenting step does not
imply a uniform bound on the total number of original rows replaced.
The stronger research target is a structural rule that finds such
repairable fibers, or an analogue of their finite-array cancellation,
inside general partial Hadamard matrices.

\subsection{A rank-eight integral completion that cannot be made a sign completion}
The residual Gram matrix for the base seed is explicitly
\[
 G=12I_{12}-P^{\mathsf T}P
   =4I_4\otimes(3I_3-J_3).
\]
It is the orthogonal sum of four scaled $A_2$ triangle Gram blocks,
with diagonal $8$ and off-diagonal $-4$ inside each triangle.
In particular, it already admits an integer factor of the correct
rank, even though it admits no sign factor.

For an explicit factor, take
\[
 C=\begin{pmatrix}
 0&1&1&1\\-1&0&1&-1\\-1&-1&0&1\\-1&1&-1&0
 \end{pmatrix},\qquad C^{\mathsf T}C=3I_4.
\]
For $j=1,\ldots,4$, let $e_j$ be a standard basis vector and $c_j$
column $j$ of $C$. Define integer vectors in $\mathbb R^8$ by
\[
 \alpha_j=\binom{2e_j}{2e_j},\qquad
 \beta_j=\binom{-e_j+c_j}{-e_j-c_j},\qquad
 \gamma_j=-\alpha_j-\beta_j.
\]
Let $Y$ have these twelve columns, grouped into four triples.
Their inner products are
\[
 \alpha_i\cdot\alpha_j=8\delta_{ij},\quad
 \beta_i\cdot\beta_j=8\delta_{ij},\quad
 \alpha_i\cdot\beta_j=-4\delta_{ij},
\]
which prove $Y^{\mathsf T}Y=G$. The nonzero eigenvalues of $G$ are
all $12$, so $YY^{\mathsf T}=12I_8$. Also $PY^{\mathsf T}=0$,
since the three columns in each block sum to zero. Thus stacking
$P$ and $Y$ gives an integer matrix $Z$ with
$ZZ^{\mathsf T}=12I_{12}$. Its completion rows have entries among
$0,\pm1,2$, not exclusively signs.

The same distinction persists at unbounded orders with exactly eight
missing rows. In the preceding tensor construction repair all but one
fiber, and call its remaining $12t-8$ sign rows $X$. If $\tau$ is the
unrepaired row of $T$, then
\[
 12tI-X^{\mathsf T}X=G\otimes(\tau^{\mathsf T}\tau).
\]
The eight integer rows $Y\otimes\tau$ give an exact orthogonal
completion, with squared norm $12t$ per row. Yet no single additional
sign row exists. Therefore integer rank-eight Gram factorization,
even with an explicit integral orthogonal completion, does not supply
the missing sign constraint. This is an explicit obstruction to a
pure integrality shortcut in residual-lattice arguments.
